% !TeX root = part01.tex

\documentclass[../final.tex]{subfiles}

\begin{document}

% ============================================================
% Практическое занятие 1
% ============================================================

\practice{1}{01.09.2026}{Классические и квантовые модели атома}


% ============================================================
% Задача 1. Модель Томсона
% ============================================================

\section*{Задача 1. Модель Томсона}


% ------------------------------------------------------------
% Условие
% ------------------------------------------------------------

    \rtxt{3}{\textbf{Условие:}}
    Согласно модели Томсона, найти радиус атома водорода и длину волны
    испускаемого им света, если энергия ионизации равна
    %
    \begin{equation*}
        I = 13.6\,\text{эВ}.
    \end{equation*}

% ------------------------------------------------------------
% Модель Томсона
% ------------------------------------------------------------

\begin{rbox}[left=18pt,right=18pt,top=14pt,bottom=10pt]{[]}
    \centering

    \includegraphics[
        width=0.42\linewidth
    ]{image/practice_01/thomson_model.png}

    \captionsetup{
        font=footnotesize,
        labelfont=bf,
        skip=5pt
    }

    \captionof{figure}{
        Модель Томсона для атома водорода.
        Положительный заряд равномерно распределён по объёму сферы
        радиуса $r_a$, электрон расположен в центре.
    }
    \label{fig:thomson_model}
\end{rbox}


В модели Томсона положительный заряд атома равномерно распределён
по объёму шара радиуса $r_a$. Поэтому плотность заряда равна
%
\begin{equation*}
    \rho
    =
    \frac{e}{\frac{4}{3}\pi r_a^3}.
\end{equation*}

Работаем в системе \rtxt{3}{СГС}, поэтому уравнение Пуассона имеет вид
%
\begin{equation*}
    \Delta\varphi(r)
    =
    -4\pi\rho(r).
\end{equation*}

Внутри атома, при $0\leq r\leq r_a$,
%
\begin{equation*}
    \Delta\varphi
    =
    -4\pi
    \frac{e}{\frac{4}{3}\pi r_a^3}
    =
    -\frac{3e}{r_a^3}.
\end{equation*}

При сферической симметрии
%
\begin{equation*}
    \Delta\varphi
    =
    \frac{1}{r^2}
    \frac{d}{dr}
    \left(
        r^2\frac{d\varphi}{dr}
    \right),
\end{equation*}
%
поэтому
%
\begin{equation*}
    \frac{1}{r^2}
    \frac{d}{dr}
    \left(
        r^2\frac{d\varphi}{dr}
    \right)
    =
    -\frac{3e}{r_a^3}.
\end{equation*}

После интегрирования
%
\begin{equation*}
    \varphi_1(r)
    =
    -\frac{er^2}{2r_a^3}
    +
    \frac{C_1}{r}
    +
    C_2.
\end{equation*}

Потенциал должен оставаться конечным при $r=0$, следовательно,
%
\begin{equation*}
    C_1=0,
\end{equation*}
%
и
%
\begin{equation*}
    \varphi_1(r)
    =
    -\frac{er^2}{2r_a^3}
    +
    C_2.
\end{equation*}

Снаружи атома, при $r\geq r_a$,
%
\begin{equation*}
    \Delta\varphi=0,
\end{equation*}
%
откуда
%
\begin{equation*}
    \varphi_2(r)
    =
    -\frac{C_3}{r}
    +
    C_4.
\end{equation*}

Так как
%
\begin{equation*}
    \varphi(\infty)=0,
\end{equation*}
%
имеем
%
\begin{equation*}
    C_4=0.
\end{equation*}


% ------------------------------------------------------------
% Сшивка решений
% ------------------------------------------------------------

\subsection*{Сшивка решений}

На границе шара потенциал непрерывен:
%
\begin{equation*}
    \varphi_1(r_a)=\varphi_2(r_a),
\end{equation*}
%
то есть
%
\begin{equation*}
    -\frac{e}{2r_a}+C_2
    =
    -\frac{C_3}{r_a}.
\end{equation*}

Кроме того, непрерывна производная потенциала:
%
\begin{equation*}
    \varphi_1'(r_a)=\varphi_2'(r_a).
\end{equation*}

Следовательно,
%
\begin{equation*}
    -\frac{e}{r_a^2}
    =
    \frac{C_3}{r_a^2},
\end{equation*}
%
откуда
%
\begin{equation*}
    C_3=-e.
\end{equation*}

Тогда
%
\begin{equation*}
    C_2=\frac{3e}{2r_a}.
\end{equation*}

Таким образом, потенциал внутри атома имеет вид

\[
    \varphi(r)
    =
    \frac{e}{r_a}
    \left(
        \frac{3}{2}
        -
        \frac{r^2}{2r_a^2}
    \right).
\]

В частности,
%
\begin{equation*}
    \varphi(0)
    =
    \frac{3e}{2r_a},
    \qquad
    \varphi(\infty)=0.
\end{equation*}


% ------------------------------------------------------------
% Энергия ионизации
% ------------------------------------------------------------

Энергия ионизации равна работе, необходимой для удаления электрона
из центра атома на бесконечность:
%
\begin{equation*}
    I
    =
    e\bigl(
        \varphi(0)-\varphi(\infty)
    \bigr).
\end{equation*}

Следовательно,
%
\begin{equation*}
    I
    =
    \frac{3e^2}{2r_a}.
\end{equation*}

Отсюда

\[
    r_a
    =
    \frac{3e^2}{2I}.
\]

Положительная ионизация:

\begin{equation*}
    I
    =
    e\bigl(
        \varphi(0)-\varphi(\infty)
    \bigr).
\end{equation*}

Численно
%
\begin{align*}
    r_a
    &=
    \frac{3}{2}
    \frac{(4.8\cdot10^{-10})^2}
         {13.6\cdot1.6\cdot10^{-12}}
    \\[3pt]
    &\approx
    1.59\cdot10^{-8}\,\text{см}
    =
    1.59\,\text{\AA}.
\end{align*}


% ============================================================
% Длина волны излучения
% ============================================================

\subsection*{Длина волны излучения}

Внутри равномерно заряженного шара
%
\begin{equation*}
    \mathbf E(r)
    =
    -\nabla\varphi
    =
    \frac{e}{r_a^3}\mathbf r.
\end{equation*}

На электрон действует сила
%
\begin{equation*}
    \mathbf F
    =
    -e\mathbf E,
\end{equation*}
%
поэтому
%
\begin{equation*}
    m\ddot{\mathbf r}
    =
    -\frac{e^2}{r_a^3}\mathbf r.
\end{equation*}

Таким образом, сила пропорциональна смещению и направлена к центру:
электрон совершает \rtxt{3}{гармонические колебания}.

Уравнение движения можно записать как
%
\begin{equation*}
    \ddot{\mathbf r}
    +
    \frac{e^2}{mr_a^3}\mathbf r
    =
    0.
\end{equation*}

Следовательно,


\[
    \omega
    =
    \sqrt{
        \frac{e^2}{mr_a^3}
    }.
\]


Длина волны испускаемого света равна
%
\begin{equation*}
    \lambda
    =
    \frac{2\pi c}{\omega}.
\end{equation*}

Поэтому

\[
    \lambda
    =
    \frac{2\pi c}{e}
    \sqrt{mr_a^3}.
\]

Подставляя
%
\begin{equation*}
    c
    =
    3\cdot10^{10}\,
    \frac{\text{см}}{\text{с}},
    \qquad
    m
    =
    9.11\cdot10^{-28}\,\text{г},
    \qquad
    r_a
    =
    1.59\cdot10^{-8}\,\text{см},
\end{equation*}
%
получаем
%
\begin{equation*}
    \lambda
    =
    \frac{2\pi\cdot3\cdot10^{10}}
         {4.8\cdot10^{-10}}
    \sqrt{
        9.11\cdot10^{-28}
        \left(
            1.59\cdot10^{-8}
        \right)^3
    }.
\end{equation*}

Отсюда
%
\begin{equation*}
    \lambda
    \approx
    23.75\cdot10^{-6}\,\text{см}
    \approx
    237.5\,\text{нм}.
\end{equation*}

\end{document}